\subsection{The complete dimension of the quadratic annihilator and an actual shift quotient}

Put $G=\widehat{GT}_0$, let $H=\ker\lambda$ be its actual
multiplier-one subgroup, and retain the literal continuous descent
homomorphism $D=D_P$ throughout.  Define
\[
 V=\{\chi\in\operatorname{Hom}_{\rm cts}(G,\mathbb F_2):
                              \chi|_{\widehat{GT}}=0\},
 \qquad \chi_n=\Delta\circ D^n\quad(n\ge0).
                                                        \tag{B.I1}
\]
We give $V$ its discrete vector-space topology and give
$\mathbb F_2^{\mathbb N_0}$ the product topology.

\begin{result}{B.2.25}{Proved}{Countably infinite quadratic annihilator and the full binary shift quotient}
The characters $\chi_n$ in (B.I1) are linearly independent,
even after restriction to $H$.  There is an actual continuous
surjective homomorphism
\[
 \begin{aligned}
 \Xi:G&\longrightarrow\mathbb F_2^{\mathbb N_0},\qquad
               g\longmapsto(\chi_0(g),\chi_1(g),\ldots),\\
 \Xi(H)&=\mathbb F_2^{\mathbb N_0},\qquad
 \Xi(\widehat{GT})=0,\qquad
 \Xi\circ D=\operatorname{shift}\circ\Xi,
 \end{aligned}                                           \tag{B.I2}
\]
where $\operatorname{shift}(a_0,a_1,\ldots)=(a_1,a_2,\ldots)$.
For every $m\ge0$, the specified finite quotient
\[
 Q_m=G\Big/\bigcap_{j=0}^{m}\ker\chi_j
       \xrightarrow[\ (\chi_0,\ldots,\chi_m)\ ]{\ \sim\ }
                                      \mathbb F_2^{m+1}       \tag{B.I3}
\]
has full gentle image and trivial GT image.  Its transition to
$Q_{m-1}$ deletes the last coordinate.  These statements give
the exact images at every level of this actual family of finite
quotients.

The full space $V$ has dimension exactly $\aleph_0$.  In
particular it is, as a discrete group, noncanonically isomorphic
to $\bigoplus_{n\ge0}\mathbb F_2$.  An exact isomorphism is
constructed in the proof by selecting a full basis extending
the displayed independent family.  If
\[
 N_V=\overline{\langle\!\langle
       [G,G],\ \{g^2:g\in G\},\ \widehat{GT}
                         \rangle\!\rangle}_G,
                                                        \tag{B.I4}
\]
then $G/N_V$ is the maximal continuous elementary abelian
two-quotient that kills $\widehat{GT}$, and evaluation on that
full basis gives a topological group isomorphism
\[
                         G/N_V\simeq\mathbb F_2^{\mathbb N_0}.
                                                        \tag{B.I5}
\]
Both $G$ and $H$ have topological generator number exactly
$\aleph_0$.
The two strict inclusions involving $\widehat{GT}$ have
the exact indices
\[
 [G:\widehat{GT}]=[\ker\Delta:\widehat{GT}]=2^{\aleph_0}.
                                                        \tag{B.I15}
\]
Here these indices count cosets; no normality of
$\widehat{GT}$ is asserted. The subgroup $\widehat{GT}$
is closed and nowhere dense in each of these two groups.
\end{result}

\begin{proof}
B.1.3 supplies an actual element $k\in H\cap\ker D$ with
$\Delta(k)=1$.  B.2.24 proves $D(H)=H$ and $D(G)=G$.
The comparison $D(F)=cFc$ on $\widehat{GT}$ from B.1.8
proves $D(\widehat{GT})=\widehat{GT}$, since $c^2=1$.
Together with A.7 this puts every $\chi_n$ in $V$.

We construct actual witnesses for all coordinate directions.
Set $g_0=k$.  Inductively choose $b_i\in H$ satisfying
$D^i(b_i)=k$, using surjectivity on $H$, and define the ordered
product
\[
              g_i=b_i\prod_{j=0}^{i-1}g_j^{\chi_j(b_i)}.
                                                        \tag{B.I6}
\]
The product is taken in increasing $j$; each exponent is the
integer zero or one represented by its displayed binary value.
We prove simultaneously that
\[
 D^i(g_i)=k,\qquad D^{i+1}(g_i)=1,\qquad
                   \chi_j(g_i)=\begin{cases}1&j=i,\\0&j\ne i.
                                  \end{cases}             \tag{B.I7}
\]
These assertions hold for $i=0$.  For $j<i$, the induction
hypothesis gives $D^i(g_j)=1$, so applying $D^i$ to (B.I6)
gives $k$.  Its next iterate is one.  Thus $\chi_i(g_i)=1$
and all later coordinates vanish.  For $\ell<i$, evaluate
the character $\chi_\ell$ on the same ordered product:
\[
 \chi_\ell(g_i)=\chi_\ell(b_i)+
       \sum_{j<i}\chi_j(b_i)\chi_\ell(g_j)
       =\chi_\ell(b_i)+\chi_\ell(b_i)=0.
                                                        \tag{B.I8}
\]
This proves (B.I7).  In particular, evaluation at $g_i$
extracts the coefficient of $\chi_i$ from every finite
linear combination of the $\chi_j$.  Independence follows
on the actual subgroup $H$ and hence on $G$.

Every coordinate of (B.I2) is a continuous homomorphism, so
$\Xi$ is a continuous homomorphism to the product.  The
elements (B.I7) map to every standard coordinate vector.
Their finite ordered products therefore give every finite
binary coordinate pattern.  Equivalently, each projection
$(\chi_0,\ldots,\chi_m):H\to\mathbb F_2^{m+1}$ is onto.
For an arbitrary infinite binary sequence, the inverse
images of its finite prefixes are nested nonempty closed
subsets of the compact group $H$.  Their intersection is
nonempty, which proves the full surjectivity in (B.I2).
Its kernel at a finite prefix is exactly the displayed
intersection in (B.I3); the first isomorphism theorem gives
that specified isomorphism and the transition map.
The GT image is zero because every coordinate belongs to
$V$.  Finally, for every $g$ and every $n$,
\[
                  \chi_n(D(g))=\chi_{n+1}(g),
                                                        \tag{B.I9}
\]
which proves the shift identity coordinate by coordinate.

We now prove the upper bound on the dimension of the entire
annihilator.  Let $F=\widehat F_2$ be the fixed free profinite
group in the definition of $G$.  For each positive integer
$r$, put
\[
 C_r=\bigcap_{\substack{U\triangleleft_o F\\{}[F:U]\le r}}U.
                                                        \tag{B.I10}
\]
There are finitely many such subgroups: a quotient of order
at most $r$ acts faithfully on at most $r$ letters, and a
homomorphism from $F$ to the corresponding symmetric group
is determined by the two generator images.  Thus $C_r$ is
open and characteristic.  The $C_r$ are nested and cofinal
among open normal subgroups, so $\bigcap_rC_r=1$.
The finite restriction maps identify
$\operatorname{Aut}(F)$ with the compatible inverse limit
of its images on $F/C_r$.  Indeed compatible automorphisms
give an automorphism of $F=\varprojlim F/C_r$, with inverse
given by the compatible inverse automorphisms.  This gives
the stated compact, metrizable automorphism topology.

The defining gentle equations cut out the closed subgroup
$G$ already used in A.4.1.  Hence the kernels of its finite
actions on $F/C_r$ form a countable identity-neighbourhood
basis.  Every continuous character of $G$ has open kernel,
contains one of these kernels, and factors through the
corresponding finite image.  Each such finite image has
only finitely many binary characters.  Therefore
$\operatorname{Hom}_{\rm cts}(G,\mathbb F_2)$, and hence
$V$, is countable.  The independent family already proved
makes its vector-space dimension exactly $\aleph_0$.

For clarity we give the isomorphism asserted for this full
space.  Order the finitely many characters of every finite
image above and concatenate these finite lists, allowing
repetitions.  Retain from this sequence precisely those
characters that vanish on $\widehat{GT}$.  Start with the
independent set $\{\chi_n:n\ge0\}$ and, in the retained
sequence order, adjoin a character exactly when it is
outside the finite-linear-combination span of the set
already chosen.  The resulting set is independent and
spans $V$: every character occurs in the sequence and is
either adjoined or already in the span at its first
occurrence.  It is countably infinite.  Enumerate it as
$(v_i)_{i\ge0}$, retaining the given $\chi_n$ as a specified
subset.  The required group isomorphism is the actual map
\[
 \bigoplus_{i\ge0}\mathbb F_2\longrightarrow V,
       (a_i)\longmapsto\sum_i a_i v_i.
                                                        \tag{B.I11}
\]
The sum is finite by its source definition; independence
and spanning prove its injectivity and surjectivity.
This is a mathematical basis selection, with no claim
that membership in the annihilator has been decided by
an effective finite test.

Every member of $V$ kills the three named families in
(B.I4), hence kills $N_V$.  Conversely $G/N_V$ is a
profinite elementary abelian two-group.  Every nonidentity
element has a nonidentity image in a finite quotient;
that quotient is an elementary abelian two-group, and a
linear functional separates its nonzero image from zero.
Pulling that functional back to $G$ gives a member of
$V$.  Consequently
\[
                         N_V=\bigcap_{v\in V}\ker v.
                                                        \tag{B.I12}
\]
Evaluation on the full basis $(v_i)$ now gives the map in
(B.I5).  It is injective by (B.I12).  Every finite coordinate
projection is onto: a proper image subspace would have a
nonzero annihilating linear functional and hence give a
nontrivial finite relation among the $v_i$.  Compactness
of the inverse images of finite prefixes proves that the
whole evaluation map is onto.  A continuous bijection
from a compact group to the Hausdorff product is a
homeomorphism, proving (B.I5) with its exact morphism.

Finally a group generated topologically by $d$ elements
has each continuous elementary abelian two-quotient of
dimension at most $d$: the generator images span that
finite quotient.  The quotients (B.I3) of unbounded
dimension, already obtained from $H$, rule out finite
topological generation for both $H$ and $G$. Both groups
are compact and metrizable by the finite-action construction
above. Choose one point in every nonempty member of a
countable topological basis. The resulting countable set
is dense and therefore generates a dense subgroup. This
proves the matching upper bound for the generator-number assertion.

For (B.I15), the actual map $\Xi$ is constant on each
right coset of $\widehat{GT}$ because its restriction to
that subgroup is zero. It therefore induces a surjection
from the coset space $G/\widehat{GT}$ onto the binary
product. On $\ker\Delta$ it induces a surjection onto
$\{0\}\times\mathbb F_2^{\mathbb N_{>0}}$, by (B.I2).
Each target has cardinality $2^{\aleph_0}$, whereas $G$
embeds in a countable product of finite groups by (B.I10)
and therefore has cardinality at most $2^{\aleph_0}$.
This proves both indices through the displayed actual
coset maps. The subgroup $\widehat{GT}$ is closed as the
zero set of the continuous pentagon word. A subgroup
with nonempty interior contains an open identity
neighbourhood and is open; an open subgroup of a compact
group has finite index. The proved indices rule this out,
so the closed subgroup has empty interior in both groups.
\end{proof}

\paragraph{The proper descent submodule and the remaining fixed-point question.}
The map
\[
 \mathbb F_2[t]\hookrightarrow V,\qquad
       \sum_n a_nt^n\longmapsto\sum_n a_n\chi_n
                                                        \tag{B.I13}
\]
is an injective module homomorphism when $t$ acts on $V$
by pullback along $D$.  The proof determines the abstract
group of all annihilating characters and this entire
explicit submodule.  B.2.27 proves that the particular
map (B.I13) is not surjective: its actual element $ak$
is killed by every $\chi_n$ and is detected by the full
character $\overline d_2$.  B.2.28 supplies the two
independent shift families generated by $\Delta$ and
$\overline d_2$, with an actual surjection onto the
product of copies of $\mathbb F_2^2$.  The full-basis
isomorphism (B.I11) remains the statement about all of
$V$; no spanning assertion for either specified
polynomial submodule is used in its proof.

On the literal fixed subgroup $C_P=\{g:D(g)=g\}$,
(B.I9) gives
\[
 \Xi(g)=(\Delta(g),\Delta(g),\ldots).
                                                        \tag{B.I14}
\]
Thus the remaining inclusion $C_P\subseteq\ker\Delta$
is precisely the assertion $\Xi(C_P)=0$.  Its possible
nonzero image is the specified diagonal subgroup of
order two.  The surjectivity statements proved above
provide preimages of its two sequences and do not
assert that those preimages are fixed by $D$.
